\chapter{Related Work}
\label{ch:literature}

This chapter reviews related research along four axes: cascade versus end-to-end
spoken dialogue systems, open speech-to-speech models and their full-duplex
capabilities, Persian speech resources, and prior work on interruption detection
and turn-taking. At the end of the chapter, the position of this project relative
to prior work is stated.

\section{From Cascades to End-to-End Models}
\label{sec:cascade-vs-e2e}

The conventional architecture of spoken dialogue systems is a chain of four
independent components: voice activity detection, speech recognition, a text
dialogue engine and text-to-speech. Its advantage is modularity: each component can
be improved, debugged and replaced with the best available model for the language
separately. For low-resource languages this modularity is a decisive advantage,
because for each component the chance of finding an acceptable Persian model is
higher than the chance of finding a Persian speech-to-speech model.

Against it, the authors of Moshi~\cite{defossez2024moshi} list three fundamental
objections to the cascade architecture. First, latency accumulates along the chain
and stretches the gap between two dialogue turns to several seconds. Second,
non-linguistic information is lost: since text is the only interface between input
and output, emotion, tone and non-speech sounds that change the meaning disappear.
Third, it depends on an explicit segmentation of speech into turns, which cannot
represent overlap, interruptions or short affirming utterances.

The observations of Chapter~\ref{ch:results} support both sides of this argument.
The latency objection shows in the measurements of the cascade track: the time to
generate a complete reply in the cascade, even on powerful hardware, is in the order
of seconds and grows with long inputs (Section~\ref{sec:res-cascade}). At the same
time the advantage of modularity also shows: only this architecture reached a
working, evaluable Persian system within the project's time frame.

\section{Open Speech-to-Speech Models}
\label{sec:openmodels}

\subsection{The Mini-Omni Family}

Mini-Omni2~\cite{xie2024miniomni2}, which the project definition names as an
example base model, is a multimodal assistant based on Qwen2-0.5B that uses Whisper
and CLIP encoders for audio and image input and generates speech with a
``text-guided delayed parallel generation'' strategy. Two of its properties matter
for our problem: first, its interruption mechanism is ``command-based'' rather than
a continuous acoustic decision; second, the project has not published a complete,
official speech-to-speech trainer, so adapting it to Persian would have required
rewriting the training pipeline from scratch.

\subsection{SpeechGPT and Discrete Speech Units}

SpeechGPT~\cite{zhang2023speechgpt} is one of the first large language models that
accepts and generates speech directly as discrete units in its vocabulary, instead
of converting speech to text through a separate recognizer and synthesizer. It was
trained in three stages: pre-training to adapt to the speech modality, cross-modal
instruction tuning, and ``chain-of-modality'' instruction tuning, in which the model
first thinks the reply as text and then turns it into speech units. The idea of
relying on intermediate text inside one integrated model is what Moshi implements
as an ``inner monologue'' synchronized with the audio. SpeechGPT, however, is
turn-based and non-streaming: it takes the whole user utterance before replying and
is not full duplex, so it was not a suitable base model for the latency and
interruption goals of this project.

\subsection{The LLaMA-Omni Family}

The first version of this family~\cite{fang2025llamaomni} used a non-autoregressive
speech decoder to reduce latency. LLaMA-Omni2~\cite{fang2025llamaomni2} is a set of
modular speech models from 0.5 to 14 billion parameters built on Qwen2.5-Instruct
models; it uses a Whisper encoder for input and an autoregressive streaming speech
decoder inspired by CosyVoice~2~\cite{du2024cosyvoice2} for output. It was trained
on only 200 thousand multi-turn dialogue samples and outperformed
GLM-4-Voice~\cite{zeng2024glm4voice}, which was trained on millions of hours of
speech.

This model was the main option at the start of the project, because its
0.5-billion-parameter size fitted the target memory budget and its stack (a Whisper
encoder plus Qwen2.5) matched the available tools. But the published repository
covers inference only. The project's first attempt to build a training pipeline for
this architecture, as Section~\ref{sec:early-attempt} describes, reached an
artifact that in fact learned to reconstruct the input rather than generate a
conversational reply; this experience was the main motivation for moving to an
official trainer.

\subsection{Freeze-Omni and the Frozen-Model Approach}

Freeze-Omni~\cite{wang2025freezeomni} takes a different strategy: the weights of
the text language model stay \emph{frozen} throughout training, and only the speech
input encoder and the speech output decoder are trained. This design has two
important advantages that are very relevant to low-resource conditions: the
intelligence of the base model stays intact and ``catastrophic forgetting'' does not
occur, and training becomes possible with much less data and compute. Full-duplex
capability is obtained by adding a state classification layer on the last layer of
the language model, which predicts whether the user has interrupted the
conversation.

In hindsight, this approach is one of the most promising options for Persian and is
suggested as future work in Chapter~\ref{ch:conclusions}; at the time of this
project's decision, an official, available trainer for a new language did not
exist.

\subsection{Qwen2.5-Omni and the ``Thinker--Talker'' Architecture}

Qwen2.5-Omni~\cite{xu2025qwen25omni} is an end-to-end multimodal model whose
``Thinker--Talker'' architecture enables simultaneous text and speech generation:
the Thinker is a language model that generates text, and the Talker is a dual-track
autoregressive model that generates audio tokens directly from the Thinker's latent
representations. It was set aside for our project for two reasons: its
7-billion-parameter version needs far more memory than the target range of
12--24 GB, and its published spoken input and output languages do not include
Persian.

\subsection{Moshi and the Final Choice of the Direct Track}

Among the models reviewed, Moshi~\cite{defossez2024moshi} was the only one that met
all three conditions this project needed: \emph{native} full-duplex capability
through parallel modelling of two audio streams, open weights, and, most
importantly, an ``official low-rank adaptation trainer'' published directly by the
model's authors~\cite{moshifinetune2024}. As explained in Section~\ref{sec:codec},
the model's ``inner monologue'' also allows the text loss and the audio loss to be
monitored separately, which was of practical value for debugging the failed
experiments.

The published experience closest to this project's direct track is
J-Moshi~\cite{ohashi2025jmoshi}, which brought Moshi to Japanese. Its authors also
found Moshi's English tokenizer inefficient for Japanese, replaced it with a
Japanese SentencePiece tokenizer with 32,000 words, and re-initialized the text
embedding tables. They then pre-trained the model on about 60,000 hours of Japanese
spoken dialogue (the 69,000-hour J-CHAT corpus) on 128 V100 GPUs for 36 hours, and
afterwards fine-tuned it on 344 hours of two-channel dialogue and 602 hours of
synthetic dialogue. In other words, adapting Moshi to a new language, even with a
suitable tokenizer, required pre-training at the scale of tens of thousands of hours
of dialogue.

Closer prior work to this idea is dGSLM~\cite{nguyen2023dgslm}, which introduced
generative modelling of spoken dialogue with two parallel streams and no text
intermediary, and showed that phenomena such as overlap and backchannels can be
reproduced in a learned way.

\subsection{PersonaPlex and Role Control in a Full-Duplex Model}

PersonaPlex~\cite{roy2026personaplex} is a full-duplex speech-to-speech model built
on the Moshi architecture that conditions the assistant's role and voice with a
combined text prompt and audio sample. Its open weights are published and keep the
native full-duplex capability. For this project, however, it had two decisive
limitations: its release is English-centred and does not support Persian, and its
documented tested hardware is outside the 12--24 GB range of requirement R6. In
addition, an official trainer to adapt it to a new language was not available at the
time of this project's decision. It was therefore considered as a comparison option
and not chosen for the direct track.

\section{Persian Speech Resources}
\label{sec:persian-resources}

For Persian, the situation of resources is asymmetric across three areas.

In \textbf{speech recognition}, usable models exist. The multilingual Common Voice
corpus~\cite{ardila2020commonvoice} has a Persian portion, and localized models
based on the Conformer architecture in NeMo~\cite{kuchaiev2019nemo} have reached
acceptable quality. The recognition service used in this project belongs to this
class.

In \textbf{text-to-speech}, the most prominent public resource is
ManaTTS~\cite{qharabagh2025manatts}: the largest public single-speaker Persian
corpus, with about 86 hours of audio at 44.1 kHz under an open licence, together with
an open processing pipeline including a sentence segmentation tool and a forced
alignment method designed for low-resource languages. The Persian voice used in this
project is derived from this resource. ParsVoice~\cite{ranjbar2025parsvoice} is a
larger corpus of Persian audiobooks with several thousand hours of speech and
hundreds of speakers, built for multi-speaker synthesis; an XTTS model fine-tuned on
it was measured in Chapter~\ref{ch:results}.

The project proposal names FarsTalk and Persian VoxCeleb as example Persian corpora;
this project's search found no published corpus under these names. The closest
published conversational resource is the Corpus of Conversational Persian
Transcripts~\cite{mohammadi2019ccpt}: transcripts of about 20 hours of informal
conversation in the Tehrani dialect, whose audio is not published and which is
therefore of no use for training a speech model.

In \textbf{full-duplex conversational data}, however, there is a serious gap. What a
full-duplex model needs to learn from is natural two-sided speech with overlap,
interruptions and backchannels. The existing Persian resources are mostly
single-speaker and read and do not contain these phenomena. Persian conversational
telephone corpora exist in research archives, but they have restrictive licences and
their scale is below the 100--200 hours targeted by the project definition. This gap
is what made the derived dataset of Chapter~\ref{ch:method} necessary.

\section{Low-Latency Streaming Inference}
\label{sec:streaming-lit}

A low-latency speech system must run each component \emph{streaming}, that is, start
the output before the whole input has arrived. A common speed measure is the
\textbf{real-time factor} (RTF): the ratio of processing time to audio duration; an
RTF below one means processing does not fall behind playback. For dialogue, however,
the RTF is not enough; what the user feels is the ``latency to the first output''.

In \textbf{speech recognition}, full attention of the encoder needs the whole
utterance. The common remedy is chunked attention: the encoder processes each chunk
with a limited past context and a short look-ahead~\cite{chen2021streamingtt}. A
cache-aware Conformer encoder~\cite{noroozi2023cacheaware} limits past and future
context during training too and keeps the activations of previous chunks, so that
step-by-step inference matches training. In the \textbf{language model}, token-by-token
generation is inherently streaming; for long dialogues, a sliding window over the
key--value cache together with keeping a few initial tokens (``attention sinks'')
maintains stability~\cite{xiao2024streamingllm}. In \textbf{speech synthesis},
clause-by-clause or chunk-by-chunk generation lets playback start before the whole
reply is generated.

This project uses all three levels, but with simpler choices
(Section~\ref{sec:cascade}): user utterances are short and the CTC recognizer
processes each utterance in about 60 ms after its end is detected; the language model
streams with a prefix cache; the synthesizer works clause by clause; and speculative
end-of-speech detection starts the work before the end is confirmed.

\section{Older Adults and Voice Assistants}
\label{sec:elderly-lit}

The voices of older adults differ acoustically from those of younger adults, and this
difference harms speech recognition: in one study, the word error rate for older
speakers was about 9 percentage points higher, and the fundamental frequency, jitter,
shimmer and harmonicity differed significantly between the two
groups~\cite{vipperla2010ageing}. The same pattern appeared on this project's human
panel: the character error rate of the deployed recognizer on older speakers was
15.6\% against 8.6\% for the others (Section~\ref{sec:reselection}).

On the user side, a three-week deployment of a smart speaker in the homes of seven
older adults with little technology experience showed that information seeking,
especially health information, was the most persistent use, and that the
reliability of spoken answers was a serious concern for them; reminders were
considered useful but were used less because of doubts about their
reliability~\cite{pradhan2020olderadults}. Other workshops with older adults listed
both the opportunities and the barriers of voice interaction~\cite{kowalski2019olderadults}.
These findings explain several design choices of this project: short, simple replies
that are honest about limitations (the elderly prompt), tolerance of longer pauses in
end-of-speech detection, slightly slower speech, and referral to a doctor or the
emergency services on risky topics. The effect of these choices on the experience of
older adults can only be measured by the human study.

\section{Interruption Detection and Turn-Taking}
\label{sec:bargein-lit}

Prior work in this area has two branches.

The \textbf{classical} branch rests on hand-crafted acoustic features and statistical
detectors. The statistical-model-based voice activity detector~\cite{sohn1999vad} is
the reference example of this branch. Interruption detection in traditional telephone
dialogue systems was in practice a voice activity detector with a conservative
threshold; the known problem of these systems was that loudspeaker playback or
ambient noise was mistaken for an interruption. A comprehensive review of turn-taking
mechanisms in dialogue systems and human--robot interaction is
in~\cite{skantze2021turntaking}, which shows that turn-taking is not merely a silence
detection problem and depends on prosodic and syntactic cues. The linguistic basis of
the topic goes back to the classical work~\cite{sacks1974turntaking}, which showed that
the gaps between turns in human conversation are remarkably short and regular.

The \textbf{learned} branch places turn-taking inside the generative model. In Moshi
and dGSLM there is no separate detector, and turn-taking is part of the learned
distribution over two parallel streams. Freeze-Omni takes a middle way: a light
classification head on the language model predicts the dialogue state.

The project definition explicitly asks for the classical branch: a module based on
energy, pitch and MFCC with a target accuracy above 80\%. This project's
implementation follows that path, with the methodological difference that instead of
manual thresholding, a gradient-boosting classifier is trained on the context-window
feature vector and its decision threshold is chosen on validation sessions only.

To label interaction events in recorded data automatically, \textbf{speaker
diarization} was used. The reference tools in this area are
pyannote.audio~\cite{bredin2020pyannote} and multi-scale diarization with dynamic
weighting~\cite{park2022msdd}; the diarization service used in this project belongs to
the second family.

To train and measure the detector under near-deployment conditions, speech is mixed
with noise and room responses. The common corpora for this are
MUSAN~\cite{snyder2015musan} (music, speech and noise),
DEMAND~\cite{thiemann2013demand} (multichannel noise of various environments) and the
OpenSLR room impulse response and noise collection~\cite{ko2017rir}.

\section{Automatic Evaluation of Dialogue Quality}
\label{sec:eval-lit}

Measuring the quality of a dialogue reply with string-matching measures is not
meaningful. The approach that has become standard in recent years is to use a large
language model as a judge. The reference work in this area~\cite{zheng2023mtbench}
showed that automatic judging with strong models correlates highly with human
preference, but the same work documented its systematic biases: position bias, length
bias and self-enhancement bias towards the output of same-family models.

This project used this method as the \emph{only} semantic measure available, but with
three restrictions: blinding and randomizing the order of the arms, deterministic
decoding, and repeating each judge call and reporting the agreement rate. More
importantly, because the judge and the candidate in this project are of the same model
family, the result is explicitly labelled a ``same-family automatic approximation''
and is not counted as a substitute for human evaluation or an independent benchmark.
For human measurement, the reference protocol is the mean opinion score based on
recommendation ITU-T P.800~\cite{itu1996p800}, which this project designed but did not
run.

\section{Position of This Project}
\label{sec:positioning}

Table~\ref{tab:positioning} summarizes the position of this project relative to prior
work.

\begin{table}[ht]
\centering
\caption{This project compared with prior work on the capabilities required by the
         project definition}
\label{tab:positioning}
\small
\begin{tabular}{L{2.9cm}ccccc}
\toprule
\textbf{Work} & \textbf{Full} & \textbf{Open} & \textbf{Official} &
\textbf{Persian} & \textbf{Memory} \\
              & \textbf{duplex} & \textbf{weights} & \textbf{trainer} &
\textbf{support} & \textbf{range} \\
\midrule
Mini-Omni2   & command-based & yes & no & no & fits \\
LLaMA-Omni2  & no & yes & no & no & fits \\
Freeze-Omni  & yes & yes & partial & no & fits \\
Qwen2.5-Omni & yes & yes & no & no & does not fit \\
Moshi        & native & yes & yes & no & does not fit \\
PersonaPlex  & native & yes & partial & no & does not fit \\
\midrule
This project, cascade & external control & yes & not needed & yes & fits \\
This project, direct  & native & yes & yes & failed & does not fit \\
\bottomrule
\end{tabular}
\end{table}

\claim{The ``memory range'' column refers to the fit with the 12--24 GB range of the
project definition; for this project's direct track it is reported from the training
memory measured on the available hardware, not from a measurement on target
hardware.}

Three points follow from this comparison.

First, none of the existing end-to-end models supports Persian; any direct track
therefore has to adapt a model, and adaptation needs an official trainer. Only Moshi
met this condition completely. PersonaPlex is also natively full duplex, but it is
English-centred and its adaptation trainer was not available for Persian.

Second, the models with native full-duplex capability do not fit within the target
memory range. This conflict between requirements R2 and R6 of the project definition
is a structural constraint, not an implementation defect.

Third, this project's distinctive contribution lies elsewhere. Prior work publishes
\emph{models}; this project builds, for one specific low-resource language, ``a data
pipeline and an evidence-first evaluation framework'' and shows which claims can and
cannot be proven with this evidence. In particular, the negative result of the direct
track, reported in Section~\ref{sec:res-moshi}, is a citable finding for later
researchers: in low-rank adaptation of a speech-to-speech model to a new language, a
reduction of the teacher-forced loss is not a sufficient criterion for checkpoint
selection.
